\begin{center}
\begin{tikzpicture}[font=\small,>=stealth]
\draw[gray!40,rounded corners] (0,0) rectangle (14.2,9);
\node[font=\small\bfseries] at (7.1,8.35) {方法一：两个直角三角形，共用同一条高};

\begin{scope}[shift={(4.4,3.2)},scale=.6]

\coordinate (L) at (-5,0); \coordinate (R) at (5,0);
\coordinate (P) at (3.75,3); \coordinate (M) at (2.4,0);
\draw[blue!75!black,thick] (L)--(P);
\draw[orange!85!black,thick] (P)--(R);
\draw[gray] (L)--(R);
\draw[magenta!80!black,very thick] (M)--(P);
\foreach \q in {L,R,P,M} \fill (\q) circle (3pt);
\node[below left] at (L) {$F_1$}; \node[below right] at (R) {$F_2$};
\node[above left] at (P) {$P$}; \node[below left] at (M) {$M$};

\coordinate (H) at (3.75,0);
\draw[dashed,gray] (P)--(H);
\draw (3.45,0)--(3.45,.3)--(3.75,.3);
\fill (H) circle(3pt);\node[below] at (H) {$H$};
\draw[blue!75!black,<->] (-5,-1.8)--(3.75,-1.8) node[midway,fill=white] {$x_0+c$};
\draw[orange!85!black,<->] (3.75,-2.7)--(5,-2.7);
\node[orange!85!black,anchor=west] at (5.4,-2.7) {$|x_0-c|$};
\node[blue!75!black,above left] at (-.6,1.5) {$r_1$};
\node[orange!85!black,right] at (5.2,1.7) {$r_2$};
\end{scope}
\node[align=left,anchor=west] at (9,5.5) {$r_1^2=(x_0+c)^2+y_0^2$\\[12pt]$r_2^2=(x_0-c)^2+y_0^2$\\[12pt]相减，公共的高消去：\\[6pt]$r_1^2-r_2^2=4cx_0$};
\node[anchor=west] at (9,2.3) {$H=(x_0,0),\quad PH=|y_0|$};
\node at (7.1,.65) {图示 $0<x_0<c$；水平长度取绝对值，平方公式对左右两支均成立。};

\end{tikzpicture}
\end{center}
